\dictchar{इ}

\bhojentry{इ}{\K{\bi}}{I}{स्वर | तद्भव}{\starsThree}%
{हिंदी तथा संस्कृत का तीसरा स्वर वर्ण; इकार।}%
{इ (Vowel I)}%
{Third vowel of Indic alphabet.}

\bhojentry{इजोत}{\K{\bi\bja\bmo\bta}}{Ijot}{हि. पुं. | [ठेठ]}{\starsThree}%
{प्रकाश, रोशनी, उजाला, अंजोर।}%
{इजोत (Ijot) / रोशनी / अंजोर}%
{Light, brightness, illumination.}
\bhojexample{भोर भइले चारो ओर इजोत पसर गइल।}{At dawn, light spread all around.}
\bhojsee{अंजोर}

\bhojentry{इजोरिया}{\K{\bi\bja\bmo\bra\bmi\bya\bmaa}}{Ijoriya}{हि. स्त्री. | [पूर्वी/मैथिली प्रभाव]}{\starsTwo}%
{चांदनी रात, शुक्ल पक्ष।}%
{अंजोरिया (Anjoriya) / इजोरिया}%
{Moonlit night.}
\bhojsee{अंजोरिया}

\bhojentry{इमली}{\K{\bi\bma\bla\bmii}}{Imli}{हि. स्त्री. | तद्भव}{\starsThree}%
{एक प्रसिद्ध खट्टा फल व उसका वृक्ष।}%
{इमली (Imli)}%
{Tamarind fruit.}

\bhojentry{इंसान}{\K{\bi\banus\bsa\bmaa\bna}}{Insaan}{हि. पुं. | अरबी आगत}{\starsThree}%
{आदमी, मनुष्य, मानुष।}%
{इंसान / आदमी}%
{Human being, man.}

\bhojsubentry{इंसानियत}{\K{\bi\banus\bsa\bmaa\bna\bmi\bya\bta}}{Insaaniyat}{हि. स्त्री. | अरबी आगत}{\starsThree}%
{मानवता, भलमनसी।}%
{इंसानियत / भलमनसी}%
{Humanity, humaneness.}

\bhojentry{इलाज}{\K{\bi\bla\bmaa\bja}}{Ilaaj}{हि. पुं. | अरबी आगत}{\starsThree}%
{चिकित्सा, दवा-दारू, उपचार।}%
{इलाज / दवा-दारू}%
{Medical treatment, cure.}
\bhojexample{अस्पताल में लइका के सही इलाज भइल।}{The child got proper treatment at the hospital.}

\bhojentry{इशारा}{\K{\bi\bsha\bmaa\bra\bmaa}}{Ishaara}{हि. पुं. | फारसी आगत}{\starsThree}%
{सैन, संकेत, इंगित।}%
{इशारा / सैन}%
{Gesture, sign, hint.}

\bhojsubentry{इशारा-काफी}{\K{\bi\bsha\bmaa\bra\bmaa}-\K{\bka\bmaa\bpha\bmii}}{Ishaara-kaafi}{हि. मुहा. | फारसी-अरबी}{\starsTwo}%
{समझदार के लिए छोटा संकेत ही बहुत है।}%
{इशारा काफी बा}%
{A hint is enough for the wise.}

\bhojentry{इनाम}{\K{\bi\bna\bmaa\bma}}{Inaam}{हि. पुं. | अरबी आगत}{\starsThree}%
{पुरस्कार, बख्शीश, पारितोषिक।}%
{इनाम / पुरस्कार}%
{Prize, reward, award.}

\bhojentry{इरादा}{\K{\bi\bra\bmaa\bda\bmaa}}{Iraada}{हि. पुं. | अरबी आगत}{\starsThree}%
{नियत, संकल्प, मनसा।}%
{इरादा / नियत}%
{Intention, resolve, purpose.}

\bhojsubentry{इरादतन}{\K{\bi\bra\bmaa\bda\bta\bna}}{Iraadatan}{हि. अव्य. | अरबी आगत}{\starsTwo}%
{जानबूझकर, जान-बूझ के।}%
{जान-बूझ के}%
{Deliberately, intentionally.}

\bhojentry{इल्जाम}{\K{\bi\bla\bja\bmaa\bma}}{Iljaam}{हि. पुं. | अरबी आगत}{\starsThree}%
{दोष, आरोप, कलंक।}%
{इल्जाम / दोस}%
{Accusation, blame.}

\bhojentry{इंतजार}{\K{\bi\banus\bta\bja\bmaa\bra}}{Intajaar}{हि. पुं. | अरबी आगत}{\starsThree}%
{बाट जोहना, अगोरा, प्रतीक्षा।}%
{इंतजार / अगोरा}%
{Wait, expectation.}
\bhojexample{हम स्टेशन पर तोहार इंतजार करब।}{I will wait for you at the station.}

\bhojentry{इंसाफ}{\K{\bi\banus\bsa\bmaa\bpha}}{Insaaf}{हि. पुं. | अरबी आगत}{\starsThree}%
{न्याय, निष्पक्षता, हक।}%
{इंसाफ / न्याय}%
{Justice, fairness.}

\bhojentry{इक्का}{\K{\bi\bka\bvir\bka\bmaa}}{Ikka}{हि. पुं. | तद्भव}{\starsThree}%
{एक घोड़े की सवारी गाड़ी; ताश का पहला पत्ता।}%
{इक्का}%
{Single-horse carriage; ace card.}

\bhojsubentry{इक्के-दुक्के}{\K{\bi\bka\bvir\bka\bme}-\K{\bda\bmu\bka\bvir\bka\bme}}{Ikke-dukke}{हि. वि. | [ठेठ]}{\starsThree}%
{विरले, एकाध, थोड़े-बहुत।}%
{इक्के-दुक्के / एकाध}%
{Sparse, one or two.}

\bhojsubentry{इक्केबल}{\K{\bi\bka\bvir\bka\bme\bba\bla}}{Ikkebal}{हि. पुं. | [ठेठ]}{\starsTwo}%
{अकेला बल, स्वावलंबन।}%
{अकेला बल / आपन भरोसा}%
{Sole strength, self-reliance.}

\bhojentry{इकरार}{\K{\bi\bka\bra\bmaa\bra}}{Ikraar}{हि. पुं. | अरबी आगत}{\starsThree}%
{स्वीकृति, वायदा, संविदा।}%
{इकरार / कबूल कइल}%
{Agreement, confession, pledge.}

\bhojentry{इतरावल}{\K{\bi\bta\bra\bmaa\bva\bla}}{Itraawal}{हि. क्रि. | [ठेठ]}{\starsTwo}%
{घमंड दिखाना, नखरा करना।}%
{इतरावल / अगराइल}%
{To act arrogant, flaunt.}

\bhojsubentry{इतराई}{\K{\bi\bta\bra\bmaa\bai}}{Itraai}{हि. स्त्री. | [ठेठ]}{\starsTwo}%
{घमंडी स्वभाव, नखरा।}%
{इतराई / नखरा}%
{Vanity, pompous behavior.}

\bhojentry{इनायत}{\K{\bi\bna\bmaa\bya\bta}}{Inaayat}{हि. स्त्री. | अरबी आगत}{\starsTwo}%
{कृपा, दया, मेहरबानी।}%
{कृपा / मेहरबानी}%
{Favor, grace, blessing.}

\bhojentry{इंजीनियर}{\K{\bi\banus\bja\bmi\bna\bmi\bya\bra}}{Injeaniyar}{हि. पुं. | अंग्रेजी आगत}{\starsThree}%
{अभियंता, तकनीकी विशेषज्ञ।}%
{इंजीनियर / अभियंता}%
{Engineer.}

\bhojentry{इंच}{\K{\bi\banus\bca}}{Inch}{हि. पुं. | अंग्रेजी आगत}{\starsThree}%
{माप की एक प्रसिद्ध इकाई।}%
{इंच (Inch)}%
{Inch unit of length.}

\bhojentry{इंजन}{\K{\bi\banus\bja\bna}}{Injan}{हि. पुं. | अंग्रेजी आगत}{\starsThree}%
{चालक मशीन, कल।}%
{इंजन (Engine)}%
{Engine, machine.}

\bhojentry{इंतजाम}{\K{\bi\banus\bta\bja\bmaa\bma}}{Intajaam}{हि. पुं. | अरबी आगत}{\starsThree}%
{प्रबंध, व्यवस्था, तैयारी।}%
{इंतजाम / प्रबंध}%
{Arrangement, organization.}

\bhojentry{इंतिहा}{\K{\bi\banus\bta\bmi\bha\bmaa}}{Intiha}{हि. स्त्री. | अरबी आगत}{\starsTwo}%
{सीमा, अंत, पराकाष्ठा।}%
{सीमा / अंत}%
{Extreme limit, peak.}

\bhojentry{इंदारा}{\K{\bi\banus\bda\bmaa\bra\bmaa}}{Indaara}{हि. पुं. | [ठेठ]}{\starsThree}%
{ईंटों का बना बड़ा पक्का कुआं।}%
{इंदारा (Indaara) / कुआँ}%
{Large brick-lined well.}

\bhojsubentry{इंदा}{\K{\bi\banus\bda\bmaa}}{Inda}{हि. पुं. | [ठेठ]}{\starsTwo}%
{कुएं का ऊपरी गोल जगत या घेरा।}%
{इंदा / जगत}%
{Well wall rim.}

\bhojentry{इंद्रदेव}{\K{\bi\banus\bda\bra\bda\bme\bva}}{Indradev}{सं. पुं. | तत्सम}{\starsTwo}%
{वर्षा के देवता, देवराज।}%
{इंद्रदेव}%
{King of gods, rain deity.}

\bhojentry{इंद्रधनुष}{\K{\baa\banus\bdha\bna\bmu\bssa}}{Indradhanush}{सं. पुं. | तत्सम}{\starsThree}%
{आकाश का सतरंगी धनुष, रामधनुष।}%
{रामधनुष / सतरंगी}%
{Rainbow.}

\bhojentry{इंद्रजाल}{\K{\bi\banus\bda\bra\bja\bmaa\bla}}{Indrajaal}{सं. पुं. | तत्सम}{\starsTwo}%
{जादू, मतिभ्रम, छल।}%
{जादू / छल}%
{Magic, illusion.}

\bhojentry{गिनती (इ-वर्ग)}{\K{\bi\bga\bna\bta\bmii}}{Ginati (I-varg)}{हि. संख्या | तद्भव}{\starsThree}%
{इ वर्ण से प्रारंभ होने वाली संख्याएँ (११, २१, ५१, ६१, ८१, ९१)।}%
{गिनती (Bhojpuri Numbers: 11, 21, 51, 61, 81, 91)}%
{Numbers starting with vowel I.}

\bhojsubentry{इग्यारह}{\K{\bi\bga\bya\bmaa\bra\bha}}{Igyaarha}{हि. संख्या | तद्भव}{\starsThree}%
{११ की संख्या।}%
{इग्यारह (11)}%
{Eleven (11).}

\bhojsubentry{इक्कीस}{\K{\bi\bka\bvir\bka\bmii\bsa}}{Ikkiis}{हि. संख्या | तद्भव}{\starsThree}%
{२१ की संख्या।}%
{इक्कीस (21)}%
{Twenty-one (21).}

\bhojsubentry{इक्यावन}{\K{\bi\bka\bya\bmaa\bva\bna}}{Ikyaavan}{हि. संख्या | तद्भव}{\starsThree}%
{५१ की संख्या।}%
{इक्यावन (51)}%
{Fifty-one (51).}

\bhojsubentry{इक्सठ}{\K{\bi\bka\bsa\bttha}}{Iksath}{हि. संख्या | तद्भव}{\starsThree}%
{६१ की संख्या।}%
{इक्सठ (61)}%
{Sixty-one (61).}

\bhojsubentry{इक्यासी}{\K{\bi\bka\bya\bmaa\bsa\bmii}}{Ikyaasi}{हि. संख्या | तद्भव}{\starsThree}%
{८१ की संख्या।}%
{इक्यासी (81)}%
{Eighty-one (81).}

\bhojsubentry{इक्यावे}{\K{\bi\bka\bya\bmaa\bva\bme}}{Ikyaave}{हि. संख्या | तद्भव}{\starsThree}%
{९१ की संख्या।}%
{इक्यावे (91)}%
{Ninety-one (91).}

\bhojentry{इक्कठा}{\K{\bi\bka\bvir\bka\bttha\bmaa}}{Ikkatha}{हि. वि. | तद्भव}{\starsThree}%
{एक स्थान पर जमा, बटोरा हुआ।}%
{एकट्ठा / जमा}%
{Gathered together, collected.}

\bhojentry{इकतारा}{\K{\bi\bka\bta\bmaa\bra\bmaa}}{Iktaara}{हि. पुं. | तद्भव}{\starsTwo}%
{एक तार का लोक संगीत वाद्य।}%
{इकतारा}%
{One-stringed musical instrument.}

\bhojentry{इकलौता}{\K{\bi\bka\bla\bmau\bta\bmaa}}{Iklauta}{हि. वि. | तद्भव}{\starsThree}%
{अकेला पुत्र या संतान।}%
{एकेगो / इकलौता}%
{Sole child, only son.}

\bhojentry{इकहरा}{\K{\bi\bka\bha\bra\bmaa}}{Ikahra}{हि. वि. | तद्भव}{\starsTwo}%
{एक परत वाला, दुबला-पतला शरीर।}%
{इकहरा देह}%
{Single-layered, lean body.}

\bhojentry{इच्छा}{\K{\bi\bca\bvir\bcha\bmaa}}{Iccha}{सं. स्त्री. | तत्सम}{\starsThree}%
{चाह, मनसा, अभिलाषा।}%
{इच्छा / चाह}%
{Desire, wish, longing.}

\bhojsubentry{इच्छित}{\K{\bi\bca\bvir\bcha\bmi\bta}}{Icchit}{सं. वि. | तत्सम}{\starsTwo}%
{मनोवांछित, चाहा हुआ।}%
{चाहल / इच्छित}%
{Desired, wished.}

\bhojentry{इजाजत}{\K{\bi\bja\bmaa\bja\bta}}{Ijaajat}{हि. स्त्री. | अरबी आगत}{\starsThree}%
{अनुमति, आज्ञा, सम्मति।}%
{इजाजत / हुक्म}%
{Permission, consent.}

\bhojentry{इज्जत}{\K{\bi\bja\bvir\bja\bta}}{Ijjat}{हि. स्त्री. | अरबी आगत}{\starsThree}%
{प्रतिष्ठा, आन, मान, इज्जत।}%
{इज्जत / मान}%
{Honor, respect, dignity.}

\bhojsubentry{इज्जतदार}{\K{\bi\bja\bvir\bja\bta\bda\bmaa\bra}}{Ijjatdaar}{हि. वि. | अरबी-फारसी}{\starsThree}%
{प्रतिष्ठित व्यक्ति, सम्मानित।}%
{इज्जतदार}%
{Honorable person.}

\bhojsubentry{इज्जत-आबरू}{\K{\bi\bja\bvir\bja\bta}-\K{\baa\bba\bra\bmuu}}{Ijjat-aabru}{हि. स्त्री. | अरबी-फारसी}{\starsThree}%
{मान-मर्यादा, शील।}%
{मान-मर्यादा}%
{Honor and reputation.}

\bhojentry{इठलावल}{\K{\bi\bttha\bla\bmaa\bva\bla}}{Ithlaawal}{हि. क्रि. | [ठेठ]}{\starsTwo}%
{नखरा करना, ऐंठना।}%
{इठलावल / नखरा कइल}%
{To show off, flirt, flaunt.}

\bhojentry{इतबार}{\K{\bi\bta\bba\bmaa\bra}}{Itbaar}{हि. पुं. | अरबी आगत}{\starsThree}%
{एतबार, भरोसा; अतवार (रविवार)।}%
{भरोसा / अतवार}%
{Trust, reliance; Sunday.}

\bhojentry{इतना}{\K{\bi\bta\bna\bmaa}}{Itna}{हि. वि. | तद्भव}{\starsThree}%
{अतना, यह परिमाण।}%
{इतना / अतना}%
{This much.}

\bhojentry{इतिहास}{\K{\bi\bta\bmi\bha\bmaa\bsa}}{Itihaas}{सं. पुं. | तत्सम}{\starsThree}%
{बीता समय, प्राचीन काल का इतिहास।}%
{इतिहास}%
{History.}

\bhojentry{इत्तिला}{\K{\bi\bta\bvir\bta\bmi\bla\bmaa}}{Ittila}{हि. स्त्री. | अरबी आगत}{\starsTwo}%
{सूचना, खबर।}%
{सूचना / खबर}%
{Information, notice.}

\bhojentry{इत्तफाक}{\K{\bi\bta\bvir\bta\bpha\bmaa\bka}}{Ittafaak}{हि. पुं. | अरबी आगत}{\starsThree}%
{संयोग, यकायक होना।}%
{अचके भइल / संयोग}%
{Coincidence, chance.}

\bhojentry{इत्र}{\K{\bi\bta\bvir\bra}}{Itr}{हि. पुं. | अरबी आगत}{\starsThree}%
{अतर, सुगंधित तेल।}%
{इत्र / अतर}%
{Perfume, attar.}

\bhojentry{इधर}{\K{\bi\bdha\bra}}{Idhar}{हि. अव्य. | तद्भव}{\starsThree}%
{एह ओर, इहवाँ।}%
{इधर / एह ओर}%
{Here, this side.}

\bhojentry{इनकार}{\K{\bi\bna\bka\bmaa\bra}}{Inkaar}{हि. पुं. | अरबी आगत}{\starsThree}%
{मनाही, ना कहना, अँगूठा देखावल।}%
{ना कइल / इनकार}%
{Refusal, denial.}

\bhojentry{इब}{\K{\bi\bba}}{Ib}{हि. अव्य. | [ठेठ]}{\starsTwo}%
{अब, इस समय।}%
{इब (Ib) / अबही}%
{Now, at present.}

\bhojentry{इबादत}{\K{\bi\bba\bmaa\bda\bta}}{Ibaadat}{हि. स्त्री. | अरबी आगत}{\starsTwo}%
{बंदगी, पूजा, वंदना।}%
{बंदगी / पूजा}%
{Worship, prayer.}

\bhojentry{इबारत}{\K{\bi\bba\bmaa\bra\bta}}{Ibaarat}{हि. स्त्री. | अरबी आगत}{\starsTwo}%
{लिखावट, पाठ।}%
{लिखावट}%
{Text, inscription.}

\bhojentry{इमारत}{\K{\bi\bma\bmaa\bra\bta}}{Imaarat}{हि. स्त्री. | अरबी आगत}{\starsThree}%
{भवन, महल, कोठा।}%
{भवन / कोठा}%
{Building, mansion.}

\bhojsubentry{इमाम}{\K{\bi\bma\bmaa\bma}}{Imaam}{हि. पुं. | अरबी आगत}{\starsTwo}%
{नमाज़ पढ़ाने वाले धार्मिक नेता।}%
{इमाम}%
{Religious leader, imam.}

\bhojsubentry{इमामबाड़ा}{\K{\bi\bma\bmaa\bma\bba\bmaa\bdda\bnukta\bmaa}}{Imaambaada}{हि. पुं. | अरबी-तद्भव}{\starsTwo}%
{मुहर्रम का ताजिया रखने का स्थान।}%
{इमामबाड़ा}%
{Shrine hall.}

\bhojentry{इमान}{\K{\bi\bma\bmaa\bna}}{Imaan}{हि. पुं. | अरबी आगत}{\starsThree}%
{धर्म, निष्ठा, नीयत।}%
{इमान / धर्म}%
{Faith, honesty, integrity.}

\bhojentry{इम्तिहान}{\K{\bi\bma\bvir\bta\bmi\bha\bmaa\bna}}{Imtihaan}{हि. पुं. | अरबी आगत}{\starsThree}%
{परीक्षा, कसौटी।}%
{परीक्षा / जाँच}%
{Examination, test.}

\bhojentry{इर्द-गिर्द}{\K{\bi\bra\bvir\bda}-\K{\bga\bmi\bra\bvir\bda}}{Ird-gird}{हि. अव्य. | फारसी}{\starsThree}%
{आसपास, चारों ओर।}%
{अगल-बगल / चारों ओर}%
{Around, nearby, vicinity.}

\bhojentry{इलम}{\K{\bi\bla\bma}}{Ilm}{हि. पुं. | अरबी आगत}{\starsTwo}%
{ज्ञान, विद्या, सीख।}%
{ज्ञान / विद्या}%
{Knowledge, learning.}

\bhojentry{इलाका}{\K{\bi\bla\bmaa\bka\bmaa}}{Ilaaka}{हि. पुं. | अरबी आगत}{\starsThree}%
{क्षेत्र, अंचल, देहात।}%
{क्षेत्र / इलाका}%
{Territory, region, area.}

\bhojentry{इलायची}{\K{\bi\bla\bmaa\bya\bca\bmii}}{Ilaayachi}{हि. स्त्री. | तद्भव}{\starsThree}%
{एक सुगंधित मसाला फल।}%
{इलायची}%
{Cardamom.}

\bhojentry{इल्लत}{\K{\bi\bla\bvir\bla\bta}}{Illat}{हि. स्त्री. | अरबी आगत}{\starsTwo}%
{बुरी आदत, लत, परेशानी।}%
{लत / कुटेव}%
{Vice, bad habit, trouble.}

\bhojentry{इल्ली}{\K{\bi\bla\bvir\bla\bmii}}{Illi}{हि. स्त्री. | [ठेठ/कृषि]}{\starsTwo}%
{फसल को नुकसान पहुँचाने वाला कीड़ा, सुंडी।}%
{सुंडी / पिल्लू}%
{Caterpillar, crop worm.}

\bhojentry{इश्क}{\K{\bi\bsha\bvir\bka}}{Ishq}{हि. पुं. | अरबी आगत}{\starsThree}%
{प्रेम, नेह, प्रीति।}%
{प्रेम / नेह}%
{Love, passion.}

\bhojentry{इश्तहार}{\K{\bi\bsha\bvir\bta\bha\bmaa\bra}}{Ishtahaar}{हि. पुं. | अरबी आगत}{\starsTwo}%
{विज्ञापन, पर्चा।}%
{विज्ञापन / पर्चा}%
{Advertisement, notice poster.}

\bhojentry{इसपताल}{\K{\bi\bsa\bpa\bta\bmaa\bla}}{Ispataal}{हि. पुं. | अंग्रेजी आगत}{\starsThree}%
{अस्पताल, औषधालय।}%
{अस्पताल / इसपताल}%
{Hospital.}

\bhojentry{इस्तरी}{\K{\bi\bsa\bta\bra\bmii}}{Istari}{हि. स्त्री. | पुर्तगाली आगत}{\starsThree}%
{कपड़े प्रेस करने का साधन।}%
{प्रेस / इस्तरी}%
{Clothes iron.}

\bhojentry{इस्तीफा}{\K{\bi\bsa\bta\bmii\bpha\bmaa}}{Istiifa}{हि. पुं. | अरबी आगत}{\starsThree}%
{त्यागपत्र।}%
{त्यागपत्र}%
{Resignation.}

\bhojentry{इस्तेमाल}{\K{\bi\bsa\bta\bme\bma\bmaa\bla}}{Istemaal}{हि. पुं. | अरबी आगत}{\starsThree}%
{उपयोग, प्रयोग, काम में लाना।}%
{प्रयोग / काम में लावल}%
{Use, application.}

\bhojentry{इहलोक}{\K{\bi\bha\bla\bmo\bka}}{Ihlog}{सं. पुं. | तत्सम}{\starsTwo}%
{यह संसार, मर्त्यलोक।}%
{इहलोक / संसार}%
{This mortal world.}

\bhojentry{इहवाँ}{\K{\bi\bha\bva\banus}}{Ihavañ}{हि. अव्य. | [ठेठ]}{\starsThree}%
{यहाँ, इस स्थान पर, ईहाँ।}%
{यहाँ / इहवाँ}%
{Here, in this place.}

\bhojentry{इहे}{\K{\bi\bha\bme}}{Ihe}{हि. सर्व. | [ठेठ]}{\starsThree}%
{यही (बात, विचार, कथन या स्थिति)।}%
{इहे (Ihe) / यही}%
{This very thing/matter/statement.}
\bhojexample{इहे हमनी के दोस्ती रहल!}{This very thing was our friendship!}

\bhojentry{इही}{\K{\bi\bha\bmii}}{Ihi}{हि. सर्व. | [ठेठ]}{\starsThree}%
{यही विशिष्ट प्राणी या व्यक्ति।}%
{इही (Ihi) / यही}%
{This specific person/being right here.}
\bhojexample{इही लइका हमार बेटा बा।}{This very boy is my son.}

\bhojentry{इहो}{\K{\bi\bha\bmo}}{Iho}{हि. सर्व. | [ठेठ]}{\starsThree}%
{यह भी।}%
{इहो (Iho) / यह भी}%
{This too.}
